\documentclass[11pt]{article}
% Shared preamble for the LSM crash-course series.
\usepackage[a4paper,margin=0.9in]{geometry}
\usepackage{amsmath,amssymb,amsthm,mathtools}
\usepackage[dvipsnames]{xcolor}
\usepackage{enumitem}
\usepackage{booktabs}
\usepackage{longtable}
\usepackage{array}
\usepackage{tcolorbox}
\usepackage{fancyhdr}
\usepackage[colorlinks=true,linkcolor=Blue]{hyperref}
\tcbuselibrary{skins,breakable}

\setlist{itemsep=3pt,topsep=4pt}
\pagestyle{fancy}\fancyhf{}
\rhead{\small Linear Statistical Models, B.Stat.\ 3rd yr}
\cfoot{\small\thepage}

\newcommand{\E}{\mathbb{E}}
\newcommand{\Var}{\operatorname{Var}}
\newcommand{\Cov}{\operatorname{Cov}}
\newcommand{\Prob}{\mathbb{P}}
\newcommand{\R}{\mathbb{R}}
\newcommand{\tr}{\operatorname{tr}}
\newcommand{\rank}{\operatorname{rank}}
\newcommand{\Cs}{\mathcal{C}}
\newcommand{\N}{\mathcal{N}}
\newcommand{\Ybar}{\overline{Y}}
\newcommand{\Xbar}{\overline{X}}
\newcommand{\xbar}{\bar{x}}
\newcommand{\iid}{\overset{\text{iid}}{\sim}}
\newcommand{\indep}{\overset{\text{indep}}{\sim}}
\newcommand{\dto}{\overset{d}{\longrightarrow}}
\newcommand{\dequal}{\overset{d}{=}}
\newcommand{\bY}{\mathbf{Y}}
\newcommand{\bX}{\mathbf{X}}
\newcommand{\bZ}{\mathbf{Z}}
\newcommand{\bW}{\mathbf{W}}
\newcommand{\bb}{\boldsymbol{\beta}}
\newcommand{\bmu}{\boldsymbol{\mu}}
\newcommand{\bep}{\boldsymbol{\varepsilon}}
\newcommand{\bal}{\boldsymbol{\alpha}}
\newcommand{\one}{\mathbf{1}}
\newcommand{\zero}{\mathbf{0}}
\newcommand{\I}{\mathbb{I}}
\newcommand{\prf}{\smallskip\noindent\textbf{Proof.}\ }
\renewcommand{\qed}{\hfill$\blacksquare$}

\newtcolorbox{idea}[1]{colback=Blue!4,colframe=Blue!55,title={#1},
  fonttitle=\bfseries,breakable}
\newtcolorbox{confusion}[1]{colback=BrickRed!5,colframe=BrickRed!60,
  title={Common confusion: #1},fonttitle=\bfseries,breakable}
\newtcolorbox{recipe}[1]{colback=ForestGreen!5,colframe=ForestGreen!55,title={#1},
  fonttitle=\bfseries,breakable}
\newtcolorbox{exam}[1]{colback=Orange!6,colframe=Orange!70,title={#1},
  fonttitle=\bfseries,breakable}

\lhead{\small One-Way ANOVA --- Crash Course}

\begin{document}

\begin{center}
{\LARGE\bfseries One-Way Analysis of Variance}\\[4pt]
{\large A Crash Course from First Principles}\\[3pt]
{\small Linear Statistical Models --- Lectures 11--12 (27 Aug -- 31 Aug)}\\[3pt]
{\small\bfseries\color{BrickRed} This is the topic the mid-semester guide is written around.}
\end{center}
\vspace{2pt}\hrule\vspace{10pt}

\noindent The fourth and most important of the four crash courses. Everything in Crash
Courses 1--3 was preparation for this. Read \S\ref{sec:numeric} (the fully worked numerical
problem) even if you read nothing else --- the exam guide says one question \textbf{will}
require a calculator.

\tableofcontents
\vspace{8pt}\hrule\vspace{10pt}

\section{Why this exists}

You have $r$ groups --- $r$ drugs, $r$ fertilisers, $r$ machines --- and a numerical response.
The question is whether the groups differ.

\begin{idea}{The concrete setting from the notes}
$r$ different drugs for reducing cholesterol. $Y_{ij}$ is the reduction in cholesterol after
one month for the $j$-th subject receiving the $i$-th drug:
\[
  i=1,\dots,r \ (\text{treatment index}),\qquad j=1,\dots,n_i\ (\text{replicate index}),
  \qquad n=\sum_i n_i .
\]
The design is \textbf{balanced} if all $n_i$ are equal; in general they need not be.
Subjects are assigned to treatments \emph{at random} --- the randomisation principle --- which
is what lets us attribute a difference in outcome to the treatment.
\end{idea}

With $r=2$ this is just the two-sample $t$-test. The whole point of ANOVA is to handle
$r>2$ \emph{without} doing $\binom{r}{2}$ separate $t$-tests (which would inflate the error
rate badly). It answers one global question first: \textbf{is there any difference at all?}

\section{The model, and the identifiability problem}

\subsection{Three equivalent ways to write it}

\begin{idea}{Cell-means form (the honest one)}
\[
  Y_{ij}=\mu_i+\varepsilon_{ij},\qquad \varepsilon_{ij}\iid \N(0,\sigma^2),
  \qquad i=1,\dots,r,\ \ j=1,\dots,n_i .
\]
Parameters: $\mu_1,\dots,\mu_r,\sigma^2$. This model is perfectly identifiable.
\end{idea}

But we want to talk about \emph{differential effects}: how much better is drug $i$ than
average? So define an overall mean $\mu$ and write $\alpha_i=\mu_i-\mu$:

\begin{idea}{Effects form (the useful one)}
\[
  \boxed{\ Y_{ij}=\mu+\alpha_i+\varepsilon_{ij}\ }
\]
Parameters: $\mu,\alpha_1,\dots,\alpha_r,\sigma^2$ --- that is $r+1$ parameters for only $r$
group means. \textbf{One too many.}
\end{idea}

\subsection{Why a constraint is forced on you}

Replace $\mu$ by $\mu-c$ and every $\alpha_i$ by $\alpha_i+c$:
\[
  Y_{ij}=(\mu-c)+(\alpha_i+c)+\varepsilon_{ij}
\]
is the \emph{same model}. The data can never distinguish the two parameter vectors. In the
language of Crash Course 2: the design matrix
\[
  \bX_{n\times(r+1)}=\begin{pmatrix}
    \one_{n_1} & \one_{n_1} & 0 & \cdots & 0\\
    \one_{n_2} & 0 & \one_{n_2} & \cdots & 0\\
    \vdots & & &\ddots& \\
    \one_{n_r} & 0&0&\cdots&\one_{n_r}
  \end{pmatrix}
  \quad\text{has}\quad \rho(\bX)=r,\ \text{not } r+1,
\]
because column $0$ is the sum of the other $r$. So $\bX^T\bX$ is singular, the model is
\textbf{not identifiable}, and $\mu$ and the $\alpha_i$ are \textbf{not estimable}
individually.

\subsection{The general form of the constraint}

Define the overall mean as a weighted average, $\mu=\sum_i w_i\mu_i$ with $\sum_i w_i=1$,
$0\le w_i\le 1$. Substituting $\mu_i=\mu+\alpha_i$:
\[
  \mu=\sum_i w_i(\mu+\alpha_i)=\mu\sum_i w_i+\sum_i w_i\alpha_i=\mu+\sum_i w_i\alpha_i
  \quad\Longrightarrow\quad \boxed{\ \sum_{i=1}^r w_i\alpha_i=0 .}
\]

\begin{recipe}{The two natural choices}
\[
  \textbf{(A)}\quad w_i=\frac1r \ \Rightarrow\ \sum_i \alpha_i = 0,
  \qquad\qquad
  \textbf{(B)}\quad w_i=\frac{n_i}{n}\ \Rightarrow\ \sum_i n_i\alpha_i = 0 .
\]
\textbf{(B) is the ``natural'' constraint and the one used in the exam guide.} It corresponds
to $\mu=\E(\Ybar_{\cdot\cdot})=\frac1n\sum_i n_i\mu_i$, the mean of the whole experiment. In a
balanced design ($n_i\equiv m$) the two coincide.
\end{recipe}

\begin{confusion}{``the constraint changes the model''}
It does not. The constraint is a \emph{labelling convention} you impose to pin down one
representative from a family of equivalent parameter vectors. Fitted values, SSE, the
$F$-statistic and every estimable quantity are completely unaffected by which constraint you
choose. Only the numerical values of $\hat\mu$ and $\hat\alpha_i$ change --- and those are
not estimable anyway.
\end{confusion}

\subsection{What \emph{is} estimable}

By the criterion of Crash Course 2, $\mathcal{N}(\bX)$ is spanned by $z=(1,-1,\dots,-1)^T$, so
$a^T\bb$ is estimable iff $a_0=\sum_{i\ge1}a_i$.

\begin{idea}{The estimable quantities of one-way ANOVA}
\begin{itemize}
  \item $\mu_i=\mu+\alpha_i$ \quad --- each group mean. \textbf{Estimable.}
  \item $\alpha_i-\alpha_{i'}=\mu_i-\mu_{i'}$ \quad --- pairwise differences. \textbf{Estimable.}
  \item $\sum_i c_i\alpha_i$ with $\sum_i c_i=0$ \quad --- a \textbf{contrast}. \textbf{Estimable}, and
        \[
          \sum_i c_i\alpha_i=\sum_i c_i(\mu_i-\mu)=\sum_i c_i\mu_i .
        \]
  \item $\mu$ alone, or $\alpha_i$ alone. \textbf{Not estimable} (without a constraint).
\end{itemize}
\end{idea}

\begin{confusion}{contrast vs.\ general linear combination}
A contrast is a linear combination of the $\mu_i$ (equivalently of the $\alpha_i$) whose
coefficients \textbf{sum to zero}. That is the whole definition, and it is exactly the
condition for estimability. $\mu_1-\mu_2$ is a contrast ($1-1=0$); $\mu_1+\mu_2$ is not;
$\mu_2-\tfrac12(\mu_1+\mu_3)$ is ($-\tfrac12+1-\tfrac12=0$).
\end{confusion}

\section{Least squares estimates}

\subsection{The direct route (fastest, use this in the exam)}

Because the cell-means form has no constraint, minimise
\[
  G(\mu_1,\dots,\mu_r)=\sum_{i=1}^r\sum_{j=1}^{n_i}(Y_{ij}-\mu_i)^2 .
\]
The terms separate by $i$, and $\sum_j (Y_{ij}-\mu_i)^2$ is minimised at the group mean. So
\[
  \boxed{\ \hat\mu_i=\Ybar_{i\cdot}=\frac1{n_i}\sum_{j=1}^{n_i}Y_{ij}\ }
\]
By Gauss--Markov, since each $\mu_i$ is estimable, $\Ybar_{i\cdot}$ is its \textbf{BLUE}.
Then under constraint (B):
\[
  \boxed{\ \hat\mu=\Ybar_{\cdot\cdot}=\frac1n\sum_i\sum_j Y_{ij},
  \qquad \hat\alpha_i=\Ybar_{i\cdot}-\Ybar_{\cdot\cdot}\ }
\]
and indeed $\sum_i n_i\hat\alpha_i=\sum_i n_i\Ybar_{i\cdot}-n\Ybar_{\cdot\cdot}=0$: the
constraint is automatically satisfied. Under the general constraint $\sum w_i\alpha_i=0$ the
BLUEs are $\hat\mu=\sum_i w_i\Ybar_{i\cdot}$ and $\hat\alpha_i=\Ybar_{i\cdot}-\hat\mu$.

\subsection{The projection route (this is what the lectures did)}

Write $\bY=\mu\one_n+\bX\bal+\bep$ where $\bX$ has the $r$ group-indicator columns. Let
$P_{\one}=\tfrac1n\one\one^T$ and $\tilde\bX=(I-P_{\one})\bX$. The constraint
$\sum n_i\alpha_i=0$ is exactly $\one^T\bX\bal=0$, so $\tilde\bX\bal=\bX\bal$, and
$P_{\bX_{\text{tot}}}=P_{\one}+P_{\tilde\bX}$ by the lemma of Crash Course 3. Then
\begin{align*}
  G(\mu,\bal)&=\|\bY-\mu\one-\tilde\bX\bal\|^2\\
  &=\underbrace{\|(I-P_{\bX_{\text{tot}}})\bY\|^2}_{\text{constant in }(\mu,\bal)}
   +\underbrace{\|P_{\one}\bY-\mu\one\|^2}_{\text{involves only }\mu}
   +\underbrace{\|P_{\tilde\bX}\bY-\bX\bal\|^2}_{\text{involves only }\bal},
\end{align*}
the three pieces being mutually orthogonal. Minimising the middle term gives
$\hat\mu=\Ybar_{\cdot\cdot}$; minimising the last gives
$\hat{\bal}=(\bX^T\bX)^{-1}\bX^TP_{\tilde\bX}\bY$. Since
$\bX^T\bX=\text{diag}(n_1,\dots,n_r)$ and
$\bX^TP_{\tilde\bX}\bY=\bX^T\bY-\Ybar_{\cdot\cdot}\bX^T\one
=\big(\sum_jY_{1j}-n_1\Ybar_{\cdot\cdot},\ \dots\big)^T$, we recover
$\hat\alpha_i=\Ybar_{i\cdot}-\Ybar_{\cdot\cdot}$. \qed

\subsection{The estimate of $\sigma^2$}

\[
  \boxed{\ \hat\sigma^2=\mathrm{MSE}=\frac{\mathrm{SSE}}{n-r}
  =\frac{\sum_{i=1}^r(n_i-1)s_i^2}{n-r},
  \qquad s_i^2=\frac1{n_i-1}\sum_{j=1}^{n_i}(Y_{ij}-\Ybar_{i\cdot})^2 . }
\]
This is the \textbf{pooled sample variance} --- a weighted average of the $r$ within-group
variances, with weights equal to their degrees of freedom. It is unbiased for $\sigma^2$.
Note $\sum_i(n_i-1)=n-r$.

\section{The variance decomposition}

\subsection{Definitions and the identity}

\begin{idea}{Sums of squares}
\begin{align*}
  \mathrm{SSTO}&=\sum_{i}\sum_j (Y_{ij}-\Ybar_{\cdot\cdot})^2 && \text{total variability}\\
  \mathrm{SSTR}&=\sum_{i=1}^r n_i(\Ybar_{i\cdot}-\Ybar_{\cdot\cdot})^2 && \text{\textbf{between} groups (BG)}\\
  \mathrm{SSE}&=\sum_i\sum_j (Y_{ij}-\Ybar_{i\cdot})^2 && \text{\textbf{within} groups (WG)}
\end{align*}
\[
  \boxed{\ \mathrm{SSTO}=\mathrm{SSTR}+\mathrm{SSE}\ }
\]
\end{idea}

\prf (algebraic) Insert and split: $Y_{ij}-\Ybar_{\cdot\cdot}=(Y_{ij}-\Ybar_{i\cdot})+(\Ybar_{i\cdot}-\Ybar_{\cdot\cdot})$.
Squaring and summing, the cross term is
\[
  2\sum_i(\Ybar_{i\cdot}-\Ybar_{\cdot\cdot})\underbrace{\sum_j(Y_{ij}-\Ybar_{i\cdot})}_{=\,0}=0 . \qquad\square
\]

\prf (geometric --- this is the lecture version) With $P_{\one}$ and $P_\bX$ the projections
onto the constant vector and onto the group-indicator space,
\[
  (I-P_{\one})\bY=\underbrace{(I-P_\bX)\bY}_{\perp}+\underbrace{(P_\bX-P_{\one})\bY}_{\perp} ,
\]
the two being orthogonal because $\Cs(\one)\subseteq\Cs(\bX)$. Pythagoras gives
$\bY^T(I-P_{\one})\bY=\bY^T(I-P_\bX)\bY+\bY^T(P_\bX-P_{\one})\bY$, which is
$\mathrm{SSTO}=\mathrm{SSE}+\mathrm{SSTR}$. \qed

\begin{idea}{The picture: two scenes}
\begin{center}
\begin{tabular}{@{}p{7.2cm}p{7.2cm}@{}}
\toprule
\textbf{Scene 1: groups clearly differ} & \textbf{Scene 2: groups indistinguishable}\\
\midrule
Points within each group are tightly bunched; the group centres are far apart. &
Points within each group are spread out; the group centres are close together.\\
$\mathrm{SSE}\ (\text{WG})\ \lll\ \mathrm{SSTR}\ (\text{BG})$ &
$\mathrm{SSTR}\ (\text{BG})$ is small relative to $\mathrm{SSE}\ (\text{WG})$\\
$F$ large $\Rightarrow$ \textbf{reject} $H_0$ & $F$ near $1$ $\Rightarrow$ \textbf{do not reject}\\
\bottomrule
\end{tabular}
\end{center}
The $F$-statistic is nothing but the formal version of ``how big is the spread
\emph{between} groups compared to the spread \emph{within} groups?''
\end{idea}

\subsection{Degrees of freedom}

\[
  \underbrace{n-1}_{\rho(I-P_{\one})}
  =\underbrace{r-1}_{\rho(P_\bX-P_{\one})}
  +\underbrace{n-r}_{\rho(I-P_\bX)} .
\]
The ranks add up to $n-1$, so \textbf{Fisher--Cochran} applies (after removing the mean) and
gives: $\mathrm{SSTR}$ and $\mathrm{SSE}$ are \emph{independent} chi-squares.

\subsection{Expected mean squares}

\begin{idea}{The table that explains why $F$ works}
\[
  \E(\mathrm{MSE})=\sigma^2, \qquad
  \E(\mathrm{MSTR})=\sigma^2+\frac{\sum_{i=1}^r n_i(\mu_i-\mu)^2}{r-1}
  =\sigma^2+\frac{\sum_i n_i\alpha_i^2}{r-1} .
\]
\end{idea}

Under $H_0:\alpha_1=\cdots=\alpha_r=0$ the second term vanishes and \emph{both} mean squares
estimate $\sigma^2$, so $F=\mathrm{MSTR}/\mathrm{MSE}$ should hover around $1$. Under $H_1$ the
numerator is inflated by $\sum n_i\alpha_i^2$, so $F$ tends to be large. \textbf{That is why
you reject for large $F$ and never for small $F$: the test is one-sided in $F$ even though the
alternative is two-sided in the $\alpha_i$.}

\section{The ANOVA table and the $F$-test}

\begin{idea}{The table --- reproduce this from memory}
\begin{center}
\begin{tabular}{@{}lcccc@{}}
\toprule
Source of variation & d.f. & SS & MS $=\frac{\text{SS}}{\text{d.f.}}$ & $\E(\text{MS})$ \\
\midrule
Between treatments & $r-1$ & SSTR & $\mathrm{MSTR}=\frac{\mathrm{SSTR}}{r-1}$
  & $\sigma^2+\frac{\sum_i n_i(\mu_i-\mu)^2}{r-1}$ \\[3pt]
Within (Error) & $n-r$ & SSE & $\mathrm{MSE}=\frac{\mathrm{SSE}}{n-r}$ & $\sigma^2$ \\[3pt]
\midrule
Total & $n-1$ & SSTO & & \\
\bottomrule
\end{tabular}
\end{center}
\[
  H_0:\alpha_1=\cdots=\alpha_r=0 \quad(\text{equivalently } \mu_1=\cdots=\mu_r),
  \qquad
  F=\frac{\mathrm{MSTR}}{\mathrm{MSE}}
   =\frac{\mathrm{SSTR}/(r-1)}{\mathrm{SSE}/(n-r)}\ \overset{H_0}{\sim}\ F_{r-1,\,n-r}.
\]
\textbf{Reject $H_0$ at level $\alpha$ if $F>F^{1-\alpha}_{r-1,\,n-r}$.}
\end{idea}

\textbf{Under $H_1$}, $F$ has a \textbf{non-central} $F_{r-1,n-r}(\lambda/2)$ distribution
with $\lambda=\frac1{\sigma^2}\sum_i n_i\alpha_i^2$; the power increases to $1$ as
$\lambda\to\infty$.

\begin{confusion}{$n-r$, not $n-r-1$}
The model $\mu+\alpha_i$ has $r+1$ parameters but the design matrix has rank $r$. The error
d.f.\ is $n-\rho(\bX)=n-r$. Every $t$ and $F$ in the question uses this. Writing $n-r-1$
poisons the whole answer.
\end{confusion}

\section{Confidence intervals and contrasts}

Everything here uses one distributional fact: $\Ybar_{i\cdot}\sim\N(\mu_i,\sigma^2/n_i)$
independently across $i$, independent of $\mathrm{MSE}$, and
$(n-r)\mathrm{MSE}/\sigma^2\sim\chi^2_{n-r}$. So any (estimable) linear combination,
standardised by its estimated standard error, is $t_{n-r}$.

\subsection{A single group mean}
\[
  \hat\mu_i=\Ybar_{i\cdot},\qquad \Var(\hat\mu_i)=\frac{\sigma^2}{n_i},
  \qquad
  \boxed{\ \Ybar_{i\cdot}\ \pm\ t^{1-\alpha/2}_{n-r}\sqrt{\frac{\mathrm{MSE}}{n_i}}\ }
\]

\subsection{A pairwise difference}
\[
  \widehat{\mu_i-\mu_{i'}}=\Ybar_{i\cdot}-\Ybar_{i'\cdot},
  \qquad
  \boxed{\ (\Ybar_{i\cdot}-\Ybar_{i'\cdot})\ \pm\ t^{1-\alpha/2}_{n-r}
  \sqrt{\mathrm{MSE}\Big(\tfrac1{n_i}+\tfrac1{n_{i'}}\Big)}\ }
\]
Note this is exactly the two-sample $t$ interval, except that $\sigma^2$ is estimated by
pooling \emph{all $r$ groups}, giving $n-r$ degrees of freedom instead of $n_i+n_{i'}-2$.
That extra precision is the payoff of the ANOVA framework.

\subsection{A general contrast --- the master formula}

\begin{recipe}{Contrast: estimate, variance, interval}
Let $\theta=\sum_{i=1}^r c_i\mu_i$ with $\sum_i c_i=0$ (equivalently $\theta=\sum_i c_i\alpha_i$).
\[
  \hat\theta=\sum_{i=1}^r c_i\Ybar_{i\cdot} \quad\text{is the \textbf{BLUE}},
  \qquad
  \Var(\hat\theta)=\sigma^2\sum_{i=1}^r\frac{c_i^2}{n_i},
\]
\[
  \boxed{\ \sum_i c_i\Ybar_{i\cdot}\ \pm\ t^{1-\alpha/2}_{n-r}\ \sqrt{\mathrm{MSE}}\ \sqrt{\sum_{i=1}^r\frac{c_i^2}{n_i}}\ }
\]
and the test of $H_0:\theta=0$ rejects when $\big|\hat\theta\big/\sqrt{\mathrm{MSE}\sum c_i^2/n_i}\big|>t^{1-\alpha/2}_{n-r}$.
\end{recipe}

\textbf{Why $\hat\theta$ is the BLUE.} $\theta$ is estimable (the coefficients sum to zero),
$\Ybar_{i\cdot}=\hat\mu_i$ is the LS estimate, so $\sum c_i\Ybar_{i\cdot}$ is
``$a^T\hat\bb$'' and the Gauss--Markov theorem applies directly.

\textbf{Why the variance is that.} The $\Ybar_{i\cdot}$ are independent with variances
$\sigma^2/n_i$, so $\Var(\sum c_i\Ybar_{i\cdot})=\sum c_i^2\sigma^2/n_i$.

\begin{confusion}{$\sqrt{\mathrm{MSE}}\cdot\sqrt{\sum c_i^2/n_i}$ --- one square root or two?}
It is a single square root of the product: $\sqrt{\mathrm{MSE}\sum c_i^2/n_i}$. The exam guide
writes it as $\sqrt{\mathrm{MSE}}\sqrt{\sum c_i^2/m_i}$, which is the same thing. What you must
\emph{not} do is forget the $1/n_i$ weights --- with unequal group sizes that is the most
common numerical error in the whole topic.
\end{confusion}

\section{The LS decomposition under the natural constraint}

The exam guide devotes a page to this argument, so here it is in full. We minimise
$G(\mu,\alpha_1,\dots,\alpha_r)=\sum_i\sum_j(Y_{ij}-\mu-\alpha_i)^2$ subject to
$\sum_i n_i\alpha_i=0$ $\ (\star)$.

Insert $\hat\mu=\Ybar_{\cdot\cdot}$ and $\hat\alpha_i=\Ybar_{i\cdot}-\Ybar_{\cdot\cdot}$
(so $\hat\mu+\hat\alpha_i=\Ybar_{i\cdot}$) and split the summand three ways:
\[
  Y_{ij}-\mu-\alpha_i=\underbrace{\big(Y_{ij}-(\hat\mu+\hat\alpha_i)\big)}_{=\,Y_{ij}-\Ybar_{i\cdot}}
  \;-\;\underbrace{(\mu-\hat\mu)}_{\text{constant in }i,j}
  \;-\;\underbrace{(\alpha_i-\hat\alpha_i)}_{\text{constant in }j} .
\]
Squaring and summing, all three cross terms vanish:
\begin{itemize}
  \item $\sum_i\sum_j(Y_{ij}-\Ybar_{i\cdot})(\mu-\hat\mu)=(\mu-\hat\mu)\sum_i\sum_j(Y_{ij}-\Ybar_{i\cdot})=0$;
  \item $\sum_i(\alpha_i-\hat\alpha_i)\sum_j(Y_{ij}-\Ybar_{i\cdot})=0$ likewise;
  \item $\sum_i n_i(\mu-\hat\mu)(\alpha_i-\hat\alpha_i)=(\mu-\hat\mu)\big(\sum_i n_i\alpha_i-\sum_i n_i\hat\alpha_i\big)=0$
        --- \textbf{and this is precisely where the constraint $(\star)$ is used}, both for
        $\alpha$ and for $\hat\alpha$.
\end{itemize}
Therefore
\[
  G(\mu,\bal)=\underbrace{\sum_i\sum_j(Y_{ij}-\Ybar_{i\cdot})^2}_{=\ \mathrm{SSE},\ \text{free of }\mu,\bal}
  \;+\;n(\mu-\hat\mu)^2\;+\;\sum_i n_i(\alpha_i-\hat\alpha_i)^2 .
\]
Each of the last two terms is a non-negative quadratic minimised at $\mu=\hat\mu$,
$\alpha_i=\hat\alpha_i$ --- and that choice satisfies $(\star)$, so it is admissible.
Setting $\mu=\hat\mu$ and all $\alpha_i=0$ instead (also admissible) and comparing gives
\[
  \sum_i\sum_j(Y_{ij}-\Ybar_{\cdot\cdot})^2=\sum_i\sum_j(Y_{ij}-\Ybar_{i\cdot})^2+\sum_i n_i\hat\alpha_i^2,
\]
i.e.\ $\mathrm{SSTO}=\mathrm{SSE}+\mathrm{SSTR}$ once more. \qed

\begin{idea}{Why this argument is worth knowing}
It proves \emph{simultaneously} that (a) the LSEs are the sample means, (b) they satisfy the
constraint automatically, and (c) the ANOVA identity holds --- all from one completed square.
If the exam asks you to ``derive the LSE under the natural identifiability constraint'', this
is the expected answer.
\end{idea}

\section{The fully worked numerical problem}\label{sec:numeric}

The guide is explicit: \emph{one numerical problem, bring a calculator.} It also tells you
the likely form: \textbf{you are given $\sum_i\sum_j Y_{ij}^2$ and the group totals
$T_i=\sum_j Y_{ij}$, and asked to find SSTO, SSTR, SSE and everything downstream.} So learn
the shortcut formulas, not the definitional ones.

\subsection{The shortcut formulas}

\begin{recipe}{The three-line computational recipe}
Let $T_i=\sum_j Y_{ij}$ (group totals), $G=\sum_i T_i$ (grand total), $n=\sum_i n_i$.
Compute the \textbf{correction factor} once:
\[
  \mathrm{CF}=\frac{G^2}{n} .
\]
Then
\[
  \boxed{\ \mathrm{SSTO}=\sum_i\sum_j Y_{ij}^2-\mathrm{CF},
  \qquad
  \mathrm{SSTR}=\sum_{i=1}^r\frac{T_i^2}{n_i}-\mathrm{CF},
  \qquad
  \mathrm{SSE}=\mathrm{SSTO}-\mathrm{SSTR}. }
\]
Never compute SSE from its definition --- get it by subtraction. It is faster and it
double-checks the identity.
\end{recipe}

\textit{Why:} $\sum\sum(Y_{ij}-\Ybar_{\cdot\cdot})^2=\sum\sum Y_{ij}^2-n\Ybar_{\cdot\cdot}^2$
and $n\Ybar_{\cdot\cdot}^2=G^2/n$; similarly
$\sum_i n_i(\Ybar_{i\cdot}-\Ybar_{\cdot\cdot})^2=\sum_i n_i\Ybar_{i\cdot}^2-n\Ybar_{\cdot\cdot}^2
=\sum_i T_i^2/n_i - G^2/n$.

\subsection{The problem}

\begin{exam}{Worked example}
Three drugs are compared for cholesterol reduction. The observed reductions (mg/dL) are

\begin{center}
\begin{tabular}{@{}lll@{}}
\toprule
Drug $A$ ($n_1=4$): & $12,\ 15,\ 11,\ 14$ \\
Drug $B$ ($n_2=5$): & $18,\ 20,\ 17,\ 19,\ 21$ \\
Drug $C$ ($n_3=3$): & $9,\ 13,\ 11$ \\
\bottomrule
\end{tabular}
\end{center}

\begin{enumerate}[label=(\alph*)]
  \item Construct the ANOVA table and test $H_0:\mu_1=\mu_2=\mu_3$ at the $5\%$ level.
  \item Give a $95\%$ confidence interval for $\mu_2$.
  \item Give a $95\%$ confidence interval for $\mu_1-\mu_3$ and interpret.
  \item Give a $95\%$ confidence interval for the contrast $\mu_2-\tfrac12(\mu_1+\mu_3)$.
  \item Report $\hat\mu$ and $\hat\alpha_i$ under the constraint $\sum n_i\alpha_i=0$.
\end{enumerate}
\emph{Given:} $F^{0.95}_{2,9}=4.26$, \quad $t^{0.975}_{9}=2.262$.
\end{exam}

\subsection{Step 1 --- the table of totals (do this first, always)}

\begin{center}
\begin{tabular}{@{}lccccc@{}}
\toprule
Group & $n_i$ & $T_i=\sum_j Y_{ij}$ & $\Ybar_{i\cdot}=T_i/n_i$ & $T_i^2/n_i$ & $\sum_j Y_{ij}^2$\\
\midrule
$A$ & $4$ & $52$ & $13$ & $52^2/4=676$ & $144{+}225{+}121{+}196=686$\\
$B$ & $5$ & $95$ & $19$ & $95^2/5=1805$ & $324{+}400{+}289{+}361{+}441=1815$\\
$C$ & $3$ & $33$ & $11$ & $33^2/3=363$ & $81{+}169{+}121=371$\\
\midrule
\textbf{Total} & $n=12$ & $G=180$ & $\Ybar_{\cdot\cdot}=15$ & $\mathbf{2844}$ & $\mathbf{2872}$\\
\bottomrule
\end{tabular}
\end{center}

\subsection{Step 2 --- the sums of squares}

\[
  \mathrm{CF}=\frac{G^2}{n}=\frac{180^2}{12}=\frac{32400}{12}=2700 .
\]
\[
  \mathrm{SSTO}=2872-2700=\mathbf{172},\qquad
  \mathrm{SSTR}=2844-2700=\mathbf{144},\qquad
  \mathrm{SSE}=172-144=\mathbf{28}.
\]

\subsection{Step 3 --- the ANOVA table}

\begin{center}
\begin{tabular}{@{}lcccc@{}}
\toprule
Source & d.f. & SS & MS & $F$\\
\midrule
Between treatments & $r-1=2$ & $144$ & $144/2=72$ & $72/3.1111=\mathbf{23.14}$\\
Error (within) & $n-r=9$ & $28$ & $28/9=3.1111$ & \\
\midrule
Total & $11$ & $172$ & & \\
\bottomrule
\end{tabular}
\end{center}

\textbf{Conclusion (a).} $F=23.14 > F^{0.95}_{2,9}=4.26$, so \textbf{reject $H_0$} at the
$5\%$ level. There is strong evidence that the three drugs differ in mean cholesterol
reduction. (In fact $p\approx0.00028$.)

\subsection{Step 4 --- the intervals}

Throughout, $\mathrm{MSE}=3.1111$, $\sqrt{\mathrm{MSE}}=1.7638$, $t^{0.975}_9=2.262$.

\medskip
\textbf{(b) CI for $\mu_2$.}
\[
  \Ybar_{2\cdot}\pm t\sqrt{\frac{\mathrm{MSE}}{n_2}}
  = 19 \pm 2.262\sqrt{\frac{3.1111}{5}}
  = 19 \pm 2.262(0.7888)
  = 19\pm 1.784
  = \boxed{(17.22,\ 20.78)} .
\]

\medskip
\textbf{(c) CI for $\mu_1-\mu_3$.} Here $c=(1,0,-1)$, so
$\sum c_i^2/n_i=\tfrac14+\tfrac13=0.58333$.
\[
  \hat\theta=13-11=2,\qquad
  \mathrm{se}=\sqrt{3.1111\times0.58333}=\sqrt{1.8148}=1.3472,
\]
\[
  2\pm 2.262(1.3472)=2\pm3.047=\boxed{(-1.05,\ 5.05)} .
\]
\textbf{Interpretation:} the interval \emph{contains zero}, so at the $5\%$ level we cannot
conclude that drugs $A$ and $C$ differ --- even though the overall $F$-test was highly
significant. The global test says \emph{somebody} differs; it does not say every pair does.

\medskip
\textbf{(d) CI for $\mu_2-\tfrac12(\mu_1+\mu_3)$.} Coefficients
$c=(-\tfrac12,\ 1,\ -\tfrac12)$; check $\sum c_i=0$. \checkmark
\[
  \sum\frac{c_i^2}{n_i}=\frac{0.25}{4}+\frac{1}{5}+\frac{0.25}{3}
  =0.0625+0.2+0.08333=0.34583 ,
\]
\[
  \hat\theta=19-\tfrac12(13+11)=19-12=7,\qquad
  \mathrm{se}=\sqrt{3.1111\times0.34583}=\sqrt{1.0759}=1.0373,
\]
\[
  7\pm2.262(1.0373)=7\pm2.347=\boxed{(4.65,\ 9.35)} .
\]
Zero is excluded, and comfortably: drug $B$ is significantly better than the average of $A$
and $C$.

\medskip
\textbf{(e) Effects.} $\hat\mu=\Ybar_{\cdot\cdot}=15$ and
\[
  \hat\alpha_1=13-15=-2,\qquad \hat\alpha_2=19-15=4,\qquad \hat\alpha_3=11-15=-4 .
\]
\textbf{Check the constraint:} $\sum n_i\hat\alpha_i=4(-2)+5(4)+3(-4)=-8+20-12=0$. \checkmark

\begin{recipe}{Calculator discipline --- do these checks, they cost seconds}
\begin{enumerate}
  \item $\sum_i T_i = G$. \quad (Here $52+95+33=180$.\ \checkmark)
  \item $\sum_i n_i = n$, and $\sum_i (\text{group } \sum Y^2) = $ overall $\sum\sum Y^2$.
  \item $\mathrm{SSTR}+\mathrm{SSE}=\mathrm{SSTO}$ and $(r-1)+(n-r)=n-1$.
  \item $\sum_i n_i\hat\alpha_i=0$.
  \item All three SS are \textbf{non-negative}. A negative SS means an arithmetic slip ---
        usually the correction factor.
  \item Keep $\mathrm{MSE}$ to $4$ decimals; rounding it to $3.1$ early shifts $F$ by about
        $1\%$ and can flip a borderline decision.
\end{enumerate}
\end{recipe}

\section{Loose ends worth five minutes}

\subsection{The two-sample $t$-test as the $r=2$ case}
With $r=2$,
\[
  T_{12}=\frac{\Ybar_{1\cdot}-\Ybar_{2\cdot}}{\hat\sigma\sqrt{\tfrac1{n_1}+\tfrac1{n_2}}}\sim t_{n-2},
  \qquad\text{and}\qquad T_{12}^2 = F \sim F_{1,\,n-2}.
\]
Under $H_1$ the non-centrality parameter is
$\dfrac{\mu_1-\mu_2}{\sigma\sqrt{1/n_1+1/n_2}}$.

\subsection{Why not just do all pairwise $t$-tests?}
With $r=5$ there are $10$ pairs; at level $0.05$ each, the chance of at least one false
rejection is far above $0.05$. The $F$-test controls one global error rate. If it rejects,
\emph{then} you look at contrasts --- ideally with a Bonferroni correction (divide $\alpha$ by
the number of comparisons; see Crash Course 1 \S6.5).

\subsection{The BLUE statement in the guide's language}
The guide records: \emph{for $\theta=\sum c_i\mu_i$ with $\sum c_i=0$, the BLUE is
$\sum c_i\Ybar_{i\cdot}$, with $\Var(\hat\theta)=\sigma^2\sum c_i^2/m_i$ and the
$100(1-\gamma)\%$ interval $\sum c_i\Ybar_{i\cdot}\pm t^{1-\gamma/2}_{n_T-r}\sqrt{\mathrm{MSE}}\sqrt{\sum c_i^2/m_i}$.}
That is \S6.3 above with $m_i$ for $n_i$ and $n_T$ for $n$. Same formula; be ready for either
notation.

\section{Practice, with answers}

\begin{enumerate}[label=\textbf{P\arabic*.}]
\item Four fertilisers, $n_i=6$ each ($n=24$). You are given $\sum\sum Y_{ij}^2=5120$,
      $G=336$, and group totals $T=(72,\ 96,\ 84,\ 84)$. Build the ANOVA table and test at
      $5\%$ given $F^{0.95}_{3,20}=3.10$.

\item In P1, give a $95\%$ CI for $\mu_2-\mu_1$ ($t^{0.975}_{20}=2.086$).

\item Show that $\mathrm{SSE}=\sum_i(n_i-1)s_i^2$, and hence that $\mathrm{MSE}$ is a weighted
      average of the within-group sample variances.

\item Prove $\E(\mathrm{MSTR})=\sigma^2+\frac{1}{r-1}\sum_i n_i\alpha_i^2$.

\item Why is $\hat\alpha_1=\Ybar_{1\cdot}-\Ybar_{\cdot\cdot}$ \emph{not} the BLUE of
      $\alpha_1$, strictly speaking? What is the precise correct statement?
\end{enumerate}

\subsection*{Answers}

\noindent\textbf{P1.} $\mathrm{CF}=336^2/24=112896/24=4704$. Check
$\sum T_i=72+96+84+84=336$ \checkmark.
$\sum T_i^2/n_i=(5184+9216+7056+7056)/6=28512/6=4752$.
\[
  \mathrm{SSTO}=5120-4704=416,\quad \mathrm{SSTR}=4752-4704=48,\quad \mathrm{SSE}=416-48=368.
\]
d.f.: $3,\ 20,\ 23$. $\mathrm{MSTR}=48/3=16$, $\mathrm{MSE}=368/20=18.4$,
$F=16/18.4=0.87$. Since $0.87<3.10$, \textbf{do not reject}: no evidence that the fertilisers
differ. (Note $F<1$ --- perfectly possible, and always a signal that the between-group
spread is smaller than chance would typically produce.)

\smallskip\noindent\textbf{P2.} $\Ybar_{2\cdot}-\Ybar_{1\cdot}=16-12=4$;
$\sum c_i^2/n_i=\tfrac16+\tfrac16=\tfrac13$;
$\mathrm{se}=\sqrt{18.4/3}=\sqrt{6.1333}=2.4766$. Interval
$4\pm2.086(2.4766)=4\pm5.166=(-1.17,\ 9.17)$ --- contains zero, consistent with the
non-significant $F$.

\smallskip\noindent\textbf{P3.} By definition
$s_i^2=\frac{1}{n_i-1}\sum_j(Y_{ij}-\Ybar_{i\cdot})^2$, so
$(n_i-1)s_i^2=\sum_j(Y_{ij}-\Ybar_{i\cdot})^2$; summing over $i$ gives $\mathrm{SSE}$. Hence
$\mathrm{MSE}=\sum_i(n_i-1)s_i^2\big/\sum_i(n_i-1)$, a weighted average with weights
$n_i-1$.

\smallskip\noindent\textbf{P4.} Write $\Ybar_{i\cdot}=\mu+\alpha_i+\bar\varepsilon_{i\cdot}$ and
$\Ybar_{\cdot\cdot}=\mu+\bar\varepsilon_{\cdot\cdot}$ (using $\sum n_i\alpha_i=0$). Then
$\Ybar_{i\cdot}-\Ybar_{\cdot\cdot}=\alpha_i+(\bar\varepsilon_{i\cdot}-\bar\varepsilon_{\cdot\cdot})$ and
\[
  \E(\mathrm{SSTR})=\sum_i n_i\Big[\alpha_i^2+\E(\bar\varepsilon_{i\cdot}-\bar\varepsilon_{\cdot\cdot})^2\Big]
  =\sum_i n_i\alpha_i^2+\sigma^2(r-1),
\]
because $\sum_i n_i\E(\bar\varepsilon_{i\cdot}-\bar\varepsilon_{\cdot\cdot})^2
=\sum_i n_i(\tfrac{1}{n_i}-\tfrac1n)\sigma^2=(r-1)\sigma^2$ (the cross terms vanish since
$\E[\bar\varepsilon_{i\cdot}\bar\varepsilon_{\cdot\cdot}]=\sigma^2/n$). Dividing by $r-1$ gives the result.

\smallskip\noindent\textbf{P5.} $\alpha_1$ is \emph{not estimable} in the unconstrained
model --- no unbiased estimator of it exists, so it has no BLUE. The correct statement is:
\emph{under the identifiability constraint $\sum_i n_i\alpha_i=0$}, the parameter
$\alpha_1=\mu_1-\sum_i (n_i/n)\mu_i$ is a contrast (its coefficients
$1-n_1/n,\ -n_2/n,\dots$ sum to zero), and \emph{as such} it is estimable with BLUE
$\Ybar_{1\cdot}-\Ybar_{\cdot\cdot}$. The constraint is what turns $\alpha_1$ into a
well-defined estimable function.

\vfill
\begin{center}\small\itshape
Companions: Simple Linear Regression; Matrix Algebra and the Gauss--Markov Framework;
Quadratic Forms, Nested Models and Testing.
\end{center}

\end{document}
